\snsection{ASPECTOS DE IMPACTO AMBIENTAL}
\justifying

Se denomina impacto ambiental a todo cambio que tiene lugar en el medio ambiente o en alguno de sus componentes, ya sea perjudicial o beneficioso, como resultado de una acción, conjunto de acciones o actividad. En este contexto, se considera que existe un impacto ambiental cuando existe una diferencia entre la situación del medio ambiente futuro producto de la realización de un proyecto, y la situación del medio ambiente futuro sin este, \parencite{chaer2020impacto}.

En el presente proyecto, impacto ambiental negativo de mayor consideración es producto de las impresiones 3D donde, en lugar de utilizar un filamento biodegradable como lo es el PLA (ácido poliláctico), se utiliza el PETG (tereftalato de polietileno glicol), él cual es un filamento que presenta mayor resistencia mecánica, durabilidad y flexibilidad, además soporta mayores temperaturas antes de comenzar a deformarse. Si bien no es biodegradable como el PLA, él mismo es reciclable al $100\%$, \parencite{pla_petg}.

El impacto ambiental respecto a la fabricación de las PCB se limita a la fabricación y transporte de los componentes desde la fábrica al proveedor, y desde el mismo a la ciudad de destino.

% No se que poner respecto al consumo energético, o sea, no se calculo mucho respecto de eso, pero bueno.
